针对神经网络处理器计算图中拓扑依赖、有限连续缓存和多执行单元并行之间的耦合，建立由拓扑序优化、驻留段分配和流水重排构成的调度方法。以题目提供的六个计算图为对象，分别评价最大缓存驻留量、额外数据搬运量和总执行周期。

对于最小驻留问题，将缓存申请延迟至首次使用之前，并在满足全部依赖后尽早释放，再比较编号顺序、有限就绪窗口贪心和前驱深度优先顺序。卷积小图的驻留峰值由 7560 降至 6984，其余五图保留生命周期收缩后的编号基准。实验表明，仅依据当前新增与释放量的局部贪心可能导致更差的全局峰值，因而采用多候选筛选而非固定使用单一规则。

对于物理缓存分配问题，以缓冲区驻留段为单位实施连续区间最佳适配，空间不足时采用下一次使用距离及搬运代价辅助选择换出对象。对受保护缓冲区造成的碎片，通过显式换出和重新装入恢复连续空间，全部代价计入评价。通过逐操作、逐地址的重放，验证容量、生命周期、依赖及使用时驻留条件。

对于性能优化问题，在额外搬运量不得增加的约束下，比较地址轮换策略，并将已有缓存方案转换为包含复用依赖的驻留图。固定驻留关系后，使用最早可开始时间与下游关键路径优先级重排执行单元指令。
在额外搬运量保持不变时，六个计算图的总执行时间降低 1.30\% 至 33.38\%。
注意力大图的执行时间由 225902 降至 150499 周期，额外搬运量均为 32852。该结果表明，在不增加缓存搬运开销的情况下，显式保留物理地址依赖并重新安排流水顺序能够改善执行效率。本文给出的结果为经约束验证的启发式可行解，不具有全局最优性保证。
